% =====================================================================
% Appendix: DASE strategies per SemBench semantic operator
% Drop into the main paper after \section{Conclusions}, e.g.
%   \appendix
%   % =====================================================================
% Appendix: DASE strategies per SemBench semantic operator
% Drop into the main paper after \section{Conclusions}, e.g.
%   \appendix
%   % =====================================================================
% Appendix: DASE strategies per SemBench semantic operator
% Drop into the main paper after \section{Conclusions}, e.g.
%   \appendix
%   % =====================================================================
% Appendix: DASE strategies per SemBench semantic operator
% Drop into the main paper after \section{Conclusions}, e.g.
%   \appendix
%   \input{appendix_sembench_strategies}
% =====================================================================

\appendix

\section{DASE Strategies per Semantic Operator}
\label{app:sembench-strategies}

A natural-language query compiles into a query plan in which a few steps
require LLM-grade reasoning---these are the \emph{semantic operators}
\textsc{sem\_filter}~(F), \textsc{sem\_join}~(J), \textsc{sem\_map}~(M),
\textsc{sem\_rank}~(R), and \textsc{sem\_classify}~(C). On SemBench
(\S\ref{sec:experiments}), DASE plays the role of a \emph{candidate
recommender} for a downstream semantic-operator engine (BigQuery in our
experiments): rather than letting the engine evaluate the LLM call on
every tuple of the input table, DASE uses the SemJI-resident embeddings
to recommend a small, high-recall set of \emph{top tuples} that the
downstream engine then evaluates with its native semantic operator. The
recommendation rule depends on the operator. The remainder of this
appendix walks through each of the five operators and describes how DASE
selects its top tuples; the reader can refer back to Table~\ref{tab:sembench}
for the per-query operator labelling and to Table~\ref{tab:results-sembench}
for the resulting cost/quality.

\subsection{\textsc{sem\_filter} (F)}

A \textsc{sem\_filter} call evaluates an NL Boolean predicate~$\ell$ on
every tuple of an input relation~$T$. DASE recommends the tuples
\emph{most likely to satisfy~$\ell$} by scoring rows with a contrastive
margin in the embedding space. From the predicate we elicit a small set
of positive paraphrases $P{=}\{p_1,\ldots\}$ and antithetical
paraphrases $N{=}\{n_1,\ldots\}$ (typically three each), embed them
once as retrieval-query vectors, and assign each row
\[
  s_i \;=\;
    \tfrac{1}{|P|}\sum_{p\in P}\!\cos(\mathbf{e}_p,\mathbf{e}_i)
    \;-\;
    \tfrac{1}{|N|}\sum_{n\in N}\!\cos(\mathbf{e}_n,\mathbf{e}_i).
\]
The contrastive form cancels the topic component shared between the
positive and negative anchors, which we found empirically more reliable
than single-anchor cosine. Two recommendation rules cover all SemBench
\textsc{sem\_filter} workloads:
\begin{itemize}
\item \emph{Bottom-$\alpha$ uncertain.} When the engine must label every
  row (\emph{Wildlife}~Q1 ``count zebra images''), DASE returns the rows
  with the smallest~$|s_i|$---the operator's hard cases---and asserts the
  cheap sign for the rest. The fraction $\alpha$ is fitted on a
  calibration split (typically 0.2--0.5).
\item \emph{Top-$K'$ shortlist.} When the predicate is paired with
  \texttt{LIMIT~$K$} (\emph{Movie}~Q1 ``five clearly positive
  reviews''), DASE returns the $K'\!\geq\!K$ rows of highest~$s_i$ and
  the engine's \texttt{AI.IF}+\texttt{LIMIT} short-circuits at the
  $K$-th survivor; this is what reduces \emph{Movie}~Q1 from~$26.3$\,s
  to~$0.5$\,s in Table~\ref{tab:results-sembench}.
\end{itemize}

\subsection{\textsc{sem\_join} (J)}

A \textsc{sem\_join} evaluates an NL match predicate over pairs
$(t_i,t_j)\!\in\!T_1\!\times\!T_2$, with naive cost
$|T_1|\!\cdot\!|T_2|$ LLM calls. DASE recommends the candidate \emph{pairs}
most likely to match by computing the cosine-similarity matrix
$S\in\mathbb{R}^{|T_1|\times|T_2|}$ over normalised embeddings of the
in-scope rows, with the diagonal masked for self-joins. Two
recommendation rules:
\begin{itemize}
\item \emph{Two-sided absolute thresholds} (\emph{Ecomm}~Q7 ``same
  brand, same category''). DASE forwards the pairs with
  $\tau_{\mathrm{lo}}{<}S_{ij}{<}\tau_{\mathrm{hi}}$, asserting matches
  for $S_{ij}\!\geq\!\tau_{\mathrm{hi}}$ and dropping
  $S_{ij}\!\leq\!\tau_{\mathrm{lo}}$ outright. Thresholds are fitted on
  a small calibration split using the symmetric tail-event geometry
  of~\S\ref{sec:semji}.
\item \emph{Top-$K$ poles}, used when the user wants extreme pairs
  (\emph{Movie}~Q5 ``ten same-sentiment review pairs''): DASE forwards
  the top-$K$ highest-similarity \emph{and} the top-$K$
  lowest-similarity pairs, the engine resolves both with one
  \texttt{AI.IF}+\texttt{LIMIT} call.
\end{itemize}
The recommended pair set is materialised as a tiny staging table; the
engine joins it twice back to the underlying relation to read the
column values and applies its native \texttt{AI.IF} on the concatenated
prompt. The crucial property is that the staging cardinality scales
with the band width, not the Cartesian product---unprefiltered baselines
time out on \emph{Movie}~Q5 where the in-scope $\mathrm{reviews}^2$
exceeds a million pairs.

\subsection{\textsc{sem\_map} (M)}

A \textsc{sem\_map} call applies a per-row NL transformation that
returns a free-form string (\emph{Ecomm}~Q3 brand extraction;
\emph{Ecomm}~Q12 colour labelling; \emph{MMQA}~Q4 genre clustering). The
output domain is open-vocabulary, so the contrastive-margin proxy used
for F/J does not apply. Instead, DASE exploits a different kind of
redundancy: rows whose embeddings are near-duplicates almost certainly
map to the \emph{same} string, so a single LLM call on one row labels
all of them.

DASE clusters the input embeddings with agglomerative clustering
(cosine metric, complete linkage, distance threshold $\tau_d{=}0.10$,
i.e.\ similarity~$\geq$~0.90) and recommends one tuple per cluster---the
medoid nearest the cluster centroid. The downstream engine runs
\texttt{AI.GENERATE} on this $K$-tuple shortlist (typically $K\!\ll\!|T|$;
on \emph{Ecomm}~Q3 this is roughly $90$ rows out of $500$), and the
returned label of each medoid is broadcast to its cluster via a join on
the synthetic cluster id---i.e.\ the broadcast itself is plain SQL,
not a special primitive. ARI degrades only when the threshold is loose
enough to fuse two genuinely distinct entities into one cluster, which
on SemBench~\emph{Ecomm} costs less than $0.01$ ARI.

\subsection{\textsc{sem\_rank} (R)}

A \textsc{sem\_rank} call asks the engine to score or order rows by a
rubric. DASE has two recommendation rules depending on whether the
output is a $K$-element ordered list or a per-row score:
\begin{itemize}
\item \emph{Top-$K$ retrieval} reduces to the F+R case: DASE returns
  the $K'$ rows of highest cheap score, the engine resolves with
  \texttt{AI.IF}+\texttt{LIMIT}.
\item \emph{Rubric scoring} (\emph{Movie}~Q9 ``score 1--5 how much the
  reviewer liked the movie'') needs a numeric label on every row, so
  DASE cannot drop confident rows without breaking the rank
  correlation. We embed the rubric anchors once (one prompt per
  rubric grade), compute, for each row, the argmax anchor (the cheap
  rubric label) and the top1${-}$top2 confidence. DASE then recommends
  the tuples in the bottom-$\alpha$ confidence tail to the engine for
  refinement via \texttt{AI.SCORE}. The two streams---cheap rubric
  scores on the confident bulk, LLM scores on the uncertain
  tail---are concatenated into the final ranking. We use a
  conservative $\alpha\!\approx\!0.7$ here because Spearman integrates
  over the score distribution: a single mis-ranked row degrades the
  global metric, so we trade more LLM calls for tail protection.
\end{itemize}

\subsection{\textsc{sem\_classify} (C)}

A \textsc{sem\_classify} call assigns each row to one of $K$ enumerated
classes (\emph{Cars}~Q10: $24$ problem categories for car complaints).
DASE embeds each class anchor prompt once and, for every row, computes
the per-anchor cosine vector; the argmax is the cheap predicted class
and the top1${-}$top2 gap is the confidence. DASE recommends the
bottom-$\alpha$ confidence tuples to the engine for refinement via
\texttt{AI.CLASSIFY}, while the engine accepts the cheap argmax for the
remaining rows. The output is the union of the two streams, and macro
$F_1$ on \emph{Cars}~Q10 stays within one point of the all-LLM baseline
at roughly half the cost (Table~\ref{tab:results-sembench}).

\subsection{Cost calibration}

All five operator strategies share a single per-row cost calibration:
DASE samples $k{=}10$ representative rows, runs them through
\texttt{AI.GENERATE\_BOOL} with the same prompt structure and modality
binding (text inline, image/audio via \texttt{ot.ref}) as the operator
under study with \texttt{thinking\_budget}=$0$, aggregates token counts
by category, prices each at the published Gemini-2.5-Flash rate, and
divides by~$k$. The result is a per-row USD figure that the cascade
multiplies by the number of tuples actually forwarded to the downstream
engine to report \texttt{cost\_usd}; \textsc{sem\_join} uses an
analogous per-pair calibration with two STRING parameters per sample.

% =====================================================================

\appendix

\section{DASE Strategies per Semantic Operator}
\label{app:sembench-strategies}

A natural-language query compiles into a query plan in which a few steps
require LLM-grade reasoning---these are the \emph{semantic operators}
\textsc{sem\_filter}~(F), \textsc{sem\_join}~(J), \textsc{sem\_map}~(M),
\textsc{sem\_rank}~(R), and \textsc{sem\_classify}~(C). On SemBench
(\S\ref{sec:experiments}), DASE plays the role of a \emph{candidate
recommender} for a downstream semantic-operator engine (BigQuery in our
experiments): rather than letting the engine evaluate the LLM call on
every tuple of the input table, DASE uses the SemJI-resident embeddings
to recommend a small, high-recall set of \emph{top tuples} that the
downstream engine then evaluates with its native semantic operator. The
recommendation rule depends on the operator. The remainder of this
appendix walks through each of the five operators and describes how DASE
selects its top tuples; the reader can refer back to Table~\ref{tab:sembench}
for the per-query operator labelling and to Table~\ref{tab:results-sembench}
for the resulting cost/quality.

\subsection{\textsc{sem\_filter} (F)}

A \textsc{sem\_filter} call evaluates an NL Boolean predicate~$\ell$ on
every tuple of an input relation~$T$. DASE recommends the tuples
\emph{most likely to satisfy~$\ell$} by scoring rows with a contrastive
margin in the embedding space. From the predicate we elicit a small set
of positive paraphrases $P{=}\{p_1,\ldots\}$ and antithetical
paraphrases $N{=}\{n_1,\ldots\}$ (typically three each), embed them
once as retrieval-query vectors, and assign each row
\[
  s_i \;=\;
    \tfrac{1}{|P|}\sum_{p\in P}\!\cos(\mathbf{e}_p,\mathbf{e}_i)
    \;-\;
    \tfrac{1}{|N|}\sum_{n\in N}\!\cos(\mathbf{e}_n,\mathbf{e}_i).
\]
The contrastive form cancels the topic component shared between the
positive and negative anchors, which we found empirically more reliable
than single-anchor cosine. Two recommendation rules cover all SemBench
\textsc{sem\_filter} workloads:
\begin{itemize}
\item \emph{Bottom-$\alpha$ uncertain.} When the engine must label every
  row (\emph{Wildlife}~Q1 ``count zebra images''), DASE returns the rows
  with the smallest~$|s_i|$---the operator's hard cases---and asserts the
  cheap sign for the rest. The fraction $\alpha$ is fitted on a
  calibration split (typically 0.2--0.5).
\item \emph{Top-$K'$ shortlist.} When the predicate is paired with
  \texttt{LIMIT~$K$} (\emph{Movie}~Q1 ``five clearly positive
  reviews''), DASE returns the $K'\!\geq\!K$ rows of highest~$s_i$ and
  the engine's \texttt{AI.IF}+\texttt{LIMIT} short-circuits at the
  $K$-th survivor; this is what reduces \emph{Movie}~Q1 from~$26.3$\,s
  to~$0.5$\,s in Table~\ref{tab:results-sembench}.
\end{itemize}

\subsection{\textsc{sem\_join} (J)}

A \textsc{sem\_join} evaluates an NL match predicate over pairs
$(t_i,t_j)\!\in\!T_1\!\times\!T_2$, with naive cost
$|T_1|\!\cdot\!|T_2|$ LLM calls. DASE recommends the candidate \emph{pairs}
most likely to match by computing the cosine-similarity matrix
$S\in\mathbb{R}^{|T_1|\times|T_2|}$ over normalised embeddings of the
in-scope rows, with the diagonal masked for self-joins. Two
recommendation rules:
\begin{itemize}
\item \emph{Two-sided absolute thresholds} (\emph{Ecomm}~Q7 ``same
  brand, same category''). DASE forwards the pairs with
  $\tau_{\mathrm{lo}}{<}S_{ij}{<}\tau_{\mathrm{hi}}$, asserting matches
  for $S_{ij}\!\geq\!\tau_{\mathrm{hi}}$ and dropping
  $S_{ij}\!\leq\!\tau_{\mathrm{lo}}$ outright. Thresholds are fitted on
  a small calibration split using the symmetric tail-event geometry
  of~\S\ref{sec:semji}.
\item \emph{Top-$K$ poles}, used when the user wants extreme pairs
  (\emph{Movie}~Q5 ``ten same-sentiment review pairs''): DASE forwards
  the top-$K$ highest-similarity \emph{and} the top-$K$
  lowest-similarity pairs, the engine resolves both with one
  \texttt{AI.IF}+\texttt{LIMIT} call.
\end{itemize}
The recommended pair set is materialised as a tiny staging table; the
engine joins it twice back to the underlying relation to read the
column values and applies its native \texttt{AI.IF} on the concatenated
prompt. The crucial property is that the staging cardinality scales
with the band width, not the Cartesian product---unprefiltered baselines
time out on \emph{Movie}~Q5 where the in-scope $\mathrm{reviews}^2$
exceeds a million pairs.

\subsection{\textsc{sem\_map} (M)}

A \textsc{sem\_map} call applies a per-row NL transformation that
returns a free-form string (\emph{Ecomm}~Q3 brand extraction;
\emph{Ecomm}~Q12 colour labelling; \emph{MMQA}~Q4 genre clustering). The
output domain is open-vocabulary, so the contrastive-margin proxy used
for F/J does not apply. Instead, DASE exploits a different kind of
redundancy: rows whose embeddings are near-duplicates almost certainly
map to the \emph{same} string, so a single LLM call on one row labels
all of them.

DASE clusters the input embeddings with agglomerative clustering
(cosine metric, complete linkage, distance threshold $\tau_d{=}0.10$,
i.e.\ similarity~$\geq$~0.90) and recommends one tuple per cluster---the
medoid nearest the cluster centroid. The downstream engine runs
\texttt{AI.GENERATE} on this $K$-tuple shortlist (typically $K\!\ll\!|T|$;
on \emph{Ecomm}~Q3 this is roughly $90$ rows out of $500$), and the
returned label of each medoid is broadcast to its cluster via a join on
the synthetic cluster id---i.e.\ the broadcast itself is plain SQL,
not a special primitive. ARI degrades only when the threshold is loose
enough to fuse two genuinely distinct entities into one cluster, which
on SemBench~\emph{Ecomm} costs less than $0.01$ ARI.

\subsection{\textsc{sem\_rank} (R)}

A \textsc{sem\_rank} call asks the engine to score or order rows by a
rubric. DASE has two recommendation rules depending on whether the
output is a $K$-element ordered list or a per-row score:
\begin{itemize}
\item \emph{Top-$K$ retrieval} reduces to the F+R case: DASE returns
  the $K'$ rows of highest cheap score, the engine resolves with
  \texttt{AI.IF}+\texttt{LIMIT}.
\item \emph{Rubric scoring} (\emph{Movie}~Q9 ``score 1--5 how much the
  reviewer liked the movie'') needs a numeric label on every row, so
  DASE cannot drop confident rows without breaking the rank
  correlation. We embed the rubric anchors once (one prompt per
  rubric grade), compute, for each row, the argmax anchor (the cheap
  rubric label) and the top1${-}$top2 confidence. DASE then recommends
  the tuples in the bottom-$\alpha$ confidence tail to the engine for
  refinement via \texttt{AI.SCORE}. The two streams---cheap rubric
  scores on the confident bulk, LLM scores on the uncertain
  tail---are concatenated into the final ranking. We use a
  conservative $\alpha\!\approx\!0.7$ here because Spearman integrates
  over the score distribution: a single mis-ranked row degrades the
  global metric, so we trade more LLM calls for tail protection.
\end{itemize}

\subsection{\textsc{sem\_classify} (C)}

A \textsc{sem\_classify} call assigns each row to one of $K$ enumerated
classes (\emph{Cars}~Q10: $24$ problem categories for car complaints).
DASE embeds each class anchor prompt once and, for every row, computes
the per-anchor cosine vector; the argmax is the cheap predicted class
and the top1${-}$top2 gap is the confidence. DASE recommends the
bottom-$\alpha$ confidence tuples to the engine for refinement via
\texttt{AI.CLASSIFY}, while the engine accepts the cheap argmax for the
remaining rows. The output is the union of the two streams, and macro
$F_1$ on \emph{Cars}~Q10 stays within one point of the all-LLM baseline
at roughly half the cost (Table~\ref{tab:results-sembench}).

\subsection{Cost calibration}

All five operator strategies share a single per-row cost calibration:
DASE samples $k{=}10$ representative rows, runs them through
\texttt{AI.GENERATE\_BOOL} with the same prompt structure and modality
binding (text inline, image/audio via \texttt{ot.ref}) as the operator
under study with \texttt{thinking\_budget}=$0$, aggregates token counts
by category, prices each at the published Gemini-2.5-Flash rate, and
divides by~$k$. The result is a per-row USD figure that the cascade
multiplies by the number of tuples actually forwarded to the downstream
engine to report \texttt{cost\_usd}; \textsc{sem\_join} uses an
analogous per-pair calibration with two STRING parameters per sample.

% =====================================================================

\appendix

\section{DASE Strategies per Semantic Operator}
\label{app:sembench-strategies}

A natural-language query compiles into a query plan in which a few steps
require LLM-grade reasoning---these are the \emph{semantic operators}
\textsc{sem\_filter}~(F), \textsc{sem\_join}~(J), \textsc{sem\_map}~(M),
\textsc{sem\_rank}~(R), and \textsc{sem\_classify}~(C). On SemBench
(\S\ref{sec:experiments}), DASE plays the role of a \emph{candidate
recommender} for a downstream semantic-operator engine (BigQuery in our
experiments): rather than letting the engine evaluate the LLM call on
every tuple of the input table, DASE uses the SemJI-resident embeddings
to recommend a small, high-recall set of \emph{top tuples} that the
downstream engine then evaluates with its native semantic operator. The
recommendation rule depends on the operator. The remainder of this
appendix walks through each of the five operators and describes how DASE
selects its top tuples; the reader can refer back to Table~\ref{tab:sembench}
for the per-query operator labelling and to Table~\ref{tab:results-sembench}
for the resulting cost/quality.

\subsection{\textsc{sem\_filter} (F)}

A \textsc{sem\_filter} call evaluates an NL Boolean predicate~$\ell$ on
every tuple of an input relation~$T$. DASE recommends the tuples
\emph{most likely to satisfy~$\ell$} by scoring rows with a contrastive
margin in the embedding space. From the predicate we elicit a small set
of positive paraphrases $P{=}\{p_1,\ldots\}$ and antithetical
paraphrases $N{=}\{n_1,\ldots\}$ (typically three each), embed them
once as retrieval-query vectors, and assign each row
\[
  s_i \;=\;
    \tfrac{1}{|P|}\sum_{p\in P}\!\cos(\mathbf{e}_p,\mathbf{e}_i)
    \;-\;
    \tfrac{1}{|N|}\sum_{n\in N}\!\cos(\mathbf{e}_n,\mathbf{e}_i).
\]
The contrastive form cancels the topic component shared between the
positive and negative anchors, which we found empirically more reliable
than single-anchor cosine. Two recommendation rules cover all SemBench
\textsc{sem\_filter} workloads:
\begin{itemize}
\item \emph{Bottom-$\alpha$ uncertain.} When the engine must label every
  row (\emph{Wildlife}~Q1 ``count zebra images''), DASE returns the rows
  with the smallest~$|s_i|$---the operator's hard cases---and asserts the
  cheap sign for the rest. The fraction $\alpha$ is fitted on a
  calibration split (typically 0.2--0.5).
\item \emph{Top-$K'$ shortlist.} When the predicate is paired with
  \texttt{LIMIT~$K$} (\emph{Movie}~Q1 ``five clearly positive
  reviews''), DASE returns the $K'\!\geq\!K$ rows of highest~$s_i$ and
  the engine's \texttt{AI.IF}+\texttt{LIMIT} short-circuits at the
  $K$-th survivor; this is what reduces \emph{Movie}~Q1 from~$26.3$\,s
  to~$0.5$\,s in Table~\ref{tab:results-sembench}.
\end{itemize}

\subsection{\textsc{sem\_join} (J)}

A \textsc{sem\_join} evaluates an NL match predicate over pairs
$(t_i,t_j)\!\in\!T_1\!\times\!T_2$, with naive cost
$|T_1|\!\cdot\!|T_2|$ LLM calls. DASE recommends the candidate \emph{pairs}
most likely to match by computing the cosine-similarity matrix
$S\in\mathbb{R}^{|T_1|\times|T_2|}$ over normalised embeddings of the
in-scope rows, with the diagonal masked for self-joins. Two
recommendation rules:
\begin{itemize}
\item \emph{Two-sided absolute thresholds} (\emph{Ecomm}~Q7 ``same
  brand, same category''). DASE forwards the pairs with
  $\tau_{\mathrm{lo}}{<}S_{ij}{<}\tau_{\mathrm{hi}}$, asserting matches
  for $S_{ij}\!\geq\!\tau_{\mathrm{hi}}$ and dropping
  $S_{ij}\!\leq\!\tau_{\mathrm{lo}}$ outright. Thresholds are fitted on
  a small calibration split using the symmetric tail-event geometry
  of~\S\ref{sec:semji}.
\item \emph{Top-$K$ poles}, used when the user wants extreme pairs
  (\emph{Movie}~Q5 ``ten same-sentiment review pairs''): DASE forwards
  the top-$K$ highest-similarity \emph{and} the top-$K$
  lowest-similarity pairs, the engine resolves both with one
  \texttt{AI.IF}+\texttt{LIMIT} call.
\end{itemize}
The recommended pair set is materialised as a tiny staging table; the
engine joins it twice back to the underlying relation to read the
column values and applies its native \texttt{AI.IF} on the concatenated
prompt. The crucial property is that the staging cardinality scales
with the band width, not the Cartesian product---unprefiltered baselines
time out on \emph{Movie}~Q5 where the in-scope $\mathrm{reviews}^2$
exceeds a million pairs.

\subsection{\textsc{sem\_map} (M)}

A \textsc{sem\_map} call applies a per-row NL transformation that
returns a free-form string (\emph{Ecomm}~Q3 brand extraction;
\emph{Ecomm}~Q12 colour labelling; \emph{MMQA}~Q4 genre clustering). The
output domain is open-vocabulary, so the contrastive-margin proxy used
for F/J does not apply. Instead, DASE exploits a different kind of
redundancy: rows whose embeddings are near-duplicates almost certainly
map to the \emph{same} string, so a single LLM call on one row labels
all of them.

DASE clusters the input embeddings with agglomerative clustering
(cosine metric, complete linkage, distance threshold $\tau_d{=}0.10$,
i.e.\ similarity~$\geq$~0.90) and recommends one tuple per cluster---the
medoid nearest the cluster centroid. The downstream engine runs
\texttt{AI.GENERATE} on this $K$-tuple shortlist (typically $K\!\ll\!|T|$;
on \emph{Ecomm}~Q3 this is roughly $90$ rows out of $500$), and the
returned label of each medoid is broadcast to its cluster via a join on
the synthetic cluster id---i.e.\ the broadcast itself is plain SQL,
not a special primitive. ARI degrades only when the threshold is loose
enough to fuse two genuinely distinct entities into one cluster, which
on SemBench~\emph{Ecomm} costs less than $0.01$ ARI.

\subsection{\textsc{sem\_rank} (R)}

A \textsc{sem\_rank} call asks the engine to score or order rows by a
rubric. DASE has two recommendation rules depending on whether the
output is a $K$-element ordered list or a per-row score:
\begin{itemize}
\item \emph{Top-$K$ retrieval} reduces to the F+R case: DASE returns
  the $K'$ rows of highest cheap score, the engine resolves with
  \texttt{AI.IF}+\texttt{LIMIT}.
\item \emph{Rubric scoring} (\emph{Movie}~Q9 ``score 1--5 how much the
  reviewer liked the movie'') needs a numeric label on every row, so
  DASE cannot drop confident rows without breaking the rank
  correlation. We embed the rubric anchors once (one prompt per
  rubric grade), compute, for each row, the argmax anchor (the cheap
  rubric label) and the top1${-}$top2 confidence. DASE then recommends
  the tuples in the bottom-$\alpha$ confidence tail to the engine for
  refinement via \texttt{AI.SCORE}. The two streams---cheap rubric
  scores on the confident bulk, LLM scores on the uncertain
  tail---are concatenated into the final ranking. We use a
  conservative $\alpha\!\approx\!0.7$ here because Spearman integrates
  over the score distribution: a single mis-ranked row degrades the
  global metric, so we trade more LLM calls for tail protection.
\end{itemize}

\subsection{\textsc{sem\_classify} (C)}

A \textsc{sem\_classify} call assigns each row to one of $K$ enumerated
classes (\emph{Cars}~Q10: $24$ problem categories for car complaints).
DASE embeds each class anchor prompt once and, for every row, computes
the per-anchor cosine vector; the argmax is the cheap predicted class
and the top1${-}$top2 gap is the confidence. DASE recommends the
bottom-$\alpha$ confidence tuples to the engine for refinement via
\texttt{AI.CLASSIFY}, while the engine accepts the cheap argmax for the
remaining rows. The output is the union of the two streams, and macro
$F_1$ on \emph{Cars}~Q10 stays within one point of the all-LLM baseline
at roughly half the cost (Table~\ref{tab:results-sembench}).

\subsection{Cost calibration}

All five operator strategies share a single per-row cost calibration:
DASE samples $k{=}10$ representative rows, runs them through
\texttt{AI.GENERATE\_BOOL} with the same prompt structure and modality
binding (text inline, image/audio via \texttt{ot.ref}) as the operator
under study with \texttt{thinking\_budget}=$0$, aggregates token counts
by category, prices each at the published Gemini-2.5-Flash rate, and
divides by~$k$. The result is a per-row USD figure that the cascade
multiplies by the number of tuples actually forwarded to the downstream
engine to report \texttt{cost\_usd}; \textsc{sem\_join} uses an
analogous per-pair calibration with two STRING parameters per sample.

% =====================================================================

\appendix

\section{DASE Strategies per Semantic Operator}
\label{app:sembench-strategies}

A natural-language query compiles into a query plan in which a few steps
require LLM-grade reasoning---these are the \emph{semantic operators}
\textsc{sem\_filter}~(F), \textsc{sem\_join}~(J), \textsc{sem\_map}~(M),
\textsc{sem\_rank}~(R), and \textsc{sem\_classify}~(C). On SemBench
(\S\ref{sec:experiments}), DASE plays the role of a \emph{candidate
recommender} for a downstream semantic-operator engine (BigQuery in our
experiments): rather than letting the engine evaluate the LLM call on
every tuple of the input table, DASE uses the SemJI-resident embeddings
to recommend a small, high-recall set of \emph{top tuples} that the
downstream engine then evaluates with its native semantic operator. The
recommendation rule depends on the operator. The remainder of this
appendix walks through each of the five operators and describes how DASE
selects its top tuples; the reader can refer back to Table~\ref{tab:sembench}
for the per-query operator labelling and to Table~\ref{tab:results-sembench}
for the resulting cost/quality.

\subsection{\textsc{sem\_filter} (F)}

A \textsc{sem\_filter} call evaluates an NL Boolean predicate~$\ell$ on
every tuple of an input relation~$T$. DASE recommends the tuples
\emph{most likely to satisfy~$\ell$} by scoring rows with a contrastive
margin in the embedding space. From the predicate we elicit a small set
of positive paraphrases $P{=}\{p_1,\ldots\}$ and antithetical
paraphrases $N{=}\{n_1,\ldots\}$ (typically three each), embed them
once as retrieval-query vectors, and assign each row
\[
  s_i \;=\;
    \tfrac{1}{|P|}\sum_{p\in P}\!\cos(\mathbf{e}_p,\mathbf{e}_i)
    \;-\;
    \tfrac{1}{|N|}\sum_{n\in N}\!\cos(\mathbf{e}_n,\mathbf{e}_i).
\]
The contrastive form cancels the topic component shared between the
positive and negative anchors, which we found empirically more reliable
than single-anchor cosine. Two recommendation rules cover all SemBench
\textsc{sem\_filter} workloads:
\begin{itemize}
\item \emph{Bottom-$\alpha$ uncertain.} When the engine must label every
  row (\emph{Wildlife}~Q1 ``count zebra images''), DASE returns the rows
  with the smallest~$|s_i|$---the operator's hard cases---and asserts the
  cheap sign for the rest. The fraction $\alpha$ is fitted on a
  calibration split (typically 0.2--0.5).
\item \emph{Top-$K'$ shortlist.} When the predicate is paired with
  \texttt{LIMIT~$K$} (\emph{Movie}~Q1 ``five clearly positive
  reviews''), DASE returns the $K'\!\geq\!K$ rows of highest~$s_i$ and
  the engine's \texttt{AI.IF}+\texttt{LIMIT} short-circuits at the
  $K$-th survivor; this is what reduces \emph{Movie}~Q1 from~$26.3$\,s
  to~$0.5$\,s in Table~\ref{tab:results-sembench}.
\end{itemize}

\subsection{\textsc{sem\_join} (J)}

A \textsc{sem\_join} evaluates an NL match predicate over pairs
$(t_i,t_j)\!\in\!T_1\!\times\!T_2$, with naive cost
$|T_1|\!\cdot\!|T_2|$ LLM calls. DASE recommends the candidate \emph{pairs}
most likely to match by computing the cosine-similarity matrix
$S\in\mathbb{R}^{|T_1|\times|T_2|}$ over normalised embeddings of the
in-scope rows, with the diagonal masked for self-joins. Two
recommendation rules:
\begin{itemize}
\item \emph{Two-sided absolute thresholds} (\emph{Ecomm}~Q7 ``same
  brand, same category''). DASE forwards the pairs with
  $\tau_{\mathrm{lo}}{<}S_{ij}{<}\tau_{\mathrm{hi}}$, asserting matches
  for $S_{ij}\!\geq\!\tau_{\mathrm{hi}}$ and dropping
  $S_{ij}\!\leq\!\tau_{\mathrm{lo}}$ outright. Thresholds are fitted on
  a small calibration split using the symmetric tail-event geometry
  of~\S\ref{sec:semji}.
\item \emph{Top-$K$ poles}, used when the user wants extreme pairs
  (\emph{Movie}~Q5 ``ten same-sentiment review pairs''): DASE forwards
  the top-$K$ highest-similarity \emph{and} the top-$K$
  lowest-similarity pairs, the engine resolves both with one
  \texttt{AI.IF}+\texttt{LIMIT} call.
\end{itemize}
The recommended pair set is materialised as a tiny staging table; the
engine joins it twice back to the underlying relation to read the
column values and applies its native \texttt{AI.IF} on the concatenated
prompt. The crucial property is that the staging cardinality scales
with the band width, not the Cartesian product---unprefiltered baselines
time out on \emph{Movie}~Q5 where the in-scope $\mathrm{reviews}^2$
exceeds a million pairs.

\subsection{\textsc{sem\_map} (M)}

A \textsc{sem\_map} call applies a per-row NL transformation that
returns a free-form string (\emph{Ecomm}~Q3 brand extraction;
\emph{Ecomm}~Q12 colour labelling; \emph{MMQA}~Q4 genre clustering). The
output domain is open-vocabulary, so the contrastive-margin proxy used
for F/J does not apply. Instead, DASE exploits a different kind of
redundancy: rows whose embeddings are near-duplicates almost certainly
map to the \emph{same} string, so a single LLM call on one row labels
all of them.

DASE clusters the input embeddings with agglomerative clustering
(cosine metric, complete linkage, distance threshold $\tau_d{=}0.10$,
i.e.\ similarity~$\geq$~0.90) and recommends one tuple per cluster---the
medoid nearest the cluster centroid. The downstream engine runs
\texttt{AI.GENERATE} on this $K$-tuple shortlist (typically $K\!\ll\!|T|$;
on \emph{Ecomm}~Q3 this is roughly $90$ rows out of $500$), and the
returned label of each medoid is broadcast to its cluster via a join on
the synthetic cluster id---i.e.\ the broadcast itself is plain SQL,
not a special primitive. ARI degrades only when the threshold is loose
enough to fuse two genuinely distinct entities into one cluster, which
on SemBench~\emph{Ecomm} costs less than $0.01$ ARI.

\subsection{\textsc{sem\_rank} (R)}

A \textsc{sem\_rank} call asks the engine to score or order rows by a
rubric. DASE has two recommendation rules depending on whether the
output is a $K$-element ordered list or a per-row score:
\begin{itemize}
\item \emph{Top-$K$ retrieval} reduces to the F+R case: DASE returns
  the $K'$ rows of highest cheap score, the engine resolves with
  \texttt{AI.IF}+\texttt{LIMIT}.
\item \emph{Rubric scoring} (\emph{Movie}~Q9 ``score 1--5 how much the
  reviewer liked the movie'') needs a numeric label on every row, so
  DASE cannot drop confident rows without breaking the rank
  correlation. We embed the rubric anchors once (one prompt per
  rubric grade), compute, for each row, the argmax anchor (the cheap
  rubric label) and the top1${-}$top2 confidence. DASE then recommends
  the tuples in the bottom-$\alpha$ confidence tail to the engine for
  refinement via \texttt{AI.SCORE}. The two streams---cheap rubric
  scores on the confident bulk, LLM scores on the uncertain
  tail---are concatenated into the final ranking. We use a
  conservative $\alpha\!\approx\!0.7$ here because Spearman integrates
  over the score distribution: a single mis-ranked row degrades the
  global metric, so we trade more LLM calls for tail protection.
\end{itemize}

\subsection{\textsc{sem\_classify} (C)}

A \textsc{sem\_classify} call assigns each row to one of $K$ enumerated
classes (\emph{Cars}~Q10: $24$ problem categories for car complaints).
DASE embeds each class anchor prompt once and, for every row, computes
the per-anchor cosine vector; the argmax is the cheap predicted class
and the top1${-}$top2 gap is the confidence. DASE recommends the
bottom-$\alpha$ confidence tuples to the engine for refinement via
\texttt{AI.CLASSIFY}, while the engine accepts the cheap argmax for the
remaining rows. The output is the union of the two streams, and macro
$F_1$ on \emph{Cars}~Q10 stays within one point of the all-LLM baseline
at roughly half the cost (Table~\ref{tab:results-sembench}).

\subsection{Cost calibration}

All five operator strategies share a single per-row cost calibration:
DASE samples $k{=}10$ representative rows, runs them through
\texttt{AI.GENERATE\_BOOL} with the same prompt structure and modality
binding (text inline, image/audio via \texttt{ot.ref}) as the operator
under study with \texttt{thinking\_budget}=$0$, aggregates token counts
by category, prices each at the published Gemini-2.5-Flash rate, and
divides by~$k$. The result is a per-row USD figure that the cascade
multiplies by the number of tuples actually forwarded to the downstream
engine to report \texttt{cost\_usd}; \textsc{sem\_join} uses an
analogous per-pair calibration with two STRING parameters per sample.
